\section{问题重述}

\subsection{背景}

休闲零食连锁门店需要根据历史销售数据安排订货、库存和促销活动。不同门店的客流基础不同，不同商品的销量规模和波动性也不同；同时，天气、节假日和门店活动可能与销售变化相关。因此，本题需要在理解历史销量结构的基础上，给出未来 7 天的销量预测，并分析商品关联和外部因素。

\subsection{需要解决的问题}

问题一要求分析不同门店和不同零食的销量情况，并分别预测未来 7 天总销量。该问题的输入是历史销售明细，输出包括门店层面、商品层面和门店--商品组合层面的预测结果。

问题二要求分析同一门店内不同零食销量之间的联系，整合同类零食，并预测未来 7 天总销量。该问题的关键是区分“商品销量同步波动”和“真实业务类别聚合”：相关性用于分析联系，附件类别字段用于正式类别预测。

问题三要求分析天气、节假日、活动日等因素与销量之间的关系。由于数据是观察数据而非随机实验，本文只讨论统计关联，不把外部变量写成严格因果影响。

问题四要求结合前三问模型和结论，预测各门店各种零食未来 7 天总销量，并与前序模型比较误差。该问题的关键是使用与未来 7 天预测一致的严格递推验证口径，避免把日滚动一步验证误写成未来 7 天一次性预测效果。

\subsection{数据说明}

本题使用附件一历史零售明细和附件二天气、节假日、促销活动数据。具体字段、日期范围、数据规模和清洗口径集中在第 \ref{sec:data_preprocessing} 节说明，字段说明另见项目数据字典 \path{DATA_DICTIONARY.md}。
